\documentclass[11pt]{article}
\usepackage[T1]{fontenc}
\usepackage{lmodern,amsmath,amssymb,amsthm,mathtools,microtype}
\usepackage[margin=1.15in]{geometry}
\usepackage[hidelinks]{hyperref}
\newtheorem{theorem}{Theorem}[section]
\newtheorem{proposition}[theorem]{Proposition}
\newtheorem{lemma}[theorem]{Lemma}
\newtheorem{definition}[theorem]{Definition}
\newtheorem{example}[theorem]{Example}
\title{Arithmetic Configuration Factorization}
\author{Raeez Lorgat}
\date{}
\hypersetup{
  pdftitle={Arithmetic Configuration Factorization},
  pdfauthor={Raeez Lorgat},
  pdfsubject={Chinese-remainder factorization and prime-tuple singular products}
}
\begin{document}
\maketitle

\begin{abstract}
For a finite set of integral shifts, the residue classes on which every
translate is invertible factor exactly over coprime squarefree moduli.  The
associated normalized local densities therefore form a positive convergent
Euler product whenever the set of shifts is admissible.  Convergence is
automatic: outside finitely many primes, the local factor differs from one by
$O_k(p^{-2})$.  The pair of shifts $\{0,2\}$ exhibits both the exceptional
prime and the resulting product.  Admissibility also supplies a finite
congruence carrier on which a divisor-incidence transform gives a generalized
Rayleigh quotient; a value above an integer $m$ forces at least $m+1$ prime
translates in the sampled interval.  Neither the carrier nor the local product
forces the quotient above a threshold.
\end{abstract}

\section{The finite product map}

Fix a nonempty finite set of distinct integers
$\mathcal H=\{h_1,\ldots,h_k\}$.  For a positive integer $q$, set
\[
 R_q=\mathbb Z/q\mathbb Z,
 \qquad
 G_q(\mathcal H)
 =\bigl\{a\in R_q:a+h_i\in R_q^\times\text{ for }1\leq i\leq k\bigr\}.
\]
If $q_1$ and $q_2$ are coprime, the Chinese-remainder ring isomorphism is
\[
 c_{q_1,q_2}:R_{q_1q_2}\longrightarrow R_{q_1}\times R_{q_2},
 \qquad
 [a]_{q_1q_2}\longmapsto([a]_{q_1},[a]_{q_2}).
\]

\begin{proposition}[Coprime configuration factorization]
For coprime positive integers $q_1,q_2$, restriction of $c_{q_1,q_2}$ gives
the bijection
\[
 \gamma_{q_1,q_2}:G_{q_1q_2}(\mathcal H)
 \xrightarrow{\ \sim\ }
 G_{q_1}(\mathcal H)\times G_{q_2}(\mathcal H).
\]
For three pairwise coprime moduli, the two iterated product maps agree after
the canonical identification of the two parenthesizations.
\end{proposition}

\begin{proof}
An element $a+h_i$ is a unit in $R_{q_1q_2}$ if and only if its two images
are units in $R_{q_1}$ and $R_{q_2}$.  The ring isomorphism therefore carries
the defining conditions for $G_{q_1q_2}(\mathcal H)$ precisely to the defining
conditions for the displayed product.  The associativity statement follows
because either iterated map sends $[a]_{q_1q_2q_3}$ to
$([a]_{q_1},[a]_{q_2},[a]_{q_3})$.
\end{proof}

For a prime $p$, define
\[
 \Omega_p(\mathcal H)=\{-h_i\bmod p:1\leq i\leq k\}\subset\mathbb F_p,
 \qquad \nu_p(\mathcal H)=|\Omega_p(\mathcal H)|.
\]
Then $G_p(\mathcal H)=\mathbb F_p\setminus\Omega_p(\mathcal H)$.  Iterating
$\gamma$ gives, for every positive squarefree $q$, a specified bijection
\[
 G_q(\mathcal H)\xrightarrow{\ \sim\ }
 \prod_{p\mid q}\bigl(\mathbb F_p\setminus\Omega_p(\mathcal H)\bigr),
 \qquad
 |G_q(\mathcal H)|=\prod_{p\mid q}(p-\nu_p(\mathcal H)).
\]

\section{The infinite product}

\begin{definition}
The set $\mathcal H$ is \emph{admissible} if
$\nu_p(\mathcal H)<p$ for every prime $p$.  Its local factor is
\[
 \sigma_p(\mathcal H)
 =\frac{1-\nu_p(\mathcal H)/p}{(1-1/p)^k}.
\]
\end{definition}

Admissibility is exactly the positivity of every local factor.  It removes
all congruence obstructions, but it asserts no distribution theorem.

\begin{theorem}[Positive singular product]
If $\mathcal H$ is admissible, then
\[
 \mathfrak S(\mathcal H)=\prod_p\sigma_p(\mathcal H)
\]
converges to a finite, strictly positive real number.  More precisely,
\[
 \sigma_p(\mathcal H)=1-\frac{k(k-1)}{2p^2}+O_k(p^{-3})
\]
outside a finite set of primes.
\end{theorem}

\begin{proof}
Put
\[
 \Delta_{\mathcal H}=\prod_{1\leq i<j\leq k}|h_i-h_j|,
\]
with the empty product equal to $1$.  If $p\nmid\Delta_{\mathcal H}$, the
reductions of the $h_i$ are distinct, so $\nu_p(\mathcal H)=k$.  Thus, for all
but finitely many primes,
\[
 \sigma_p=(1-k/p)(1-1/p)^{-k}.
\]
For $p>k$, the convergent power series for the logarithm gives
\[
 \log\sigma_p
 =-\sum_{n\geq2}\frac{k^n-k}{n p^n}
 =-\frac{k(k-1)}{2p^2}+O_k(p^{-3}).
\]
Exponentiation gives the asserted expansion and, in particular,
$|\sigma_p-1|=O_k(p^{-2})$.  Hence
$\sum_p|\sigma_p-1|<\infty$.  Eventually $|\sigma_p-1|\leq1/2$, and then
$|\log\sigma_p|\leq2|\sigma_p-1|$.  The logarithmic series therefore
converges absolutely.  All its factors are positive by admissibility, so the
product is finite and strictly positive.
\end{proof}

\begin{example}[The shifts $\{0,2\}$]
Let $\mathcal H=\{0,2\}$.  At the exceptional prime $2$, the two forbidden
classes coincide:
\[
 \Omega_2(\mathcal H)=\{0\},\qquad
 G_2(\mathcal H)=\{1\},\qquad
 \sigma_2=\frac{1-1/2}{(1-1/2)^2}=2.
\]
For every odd prime $p$, the forbidden classes $0$ and $-2$ are distinct, so
\[
 |G_p(\mathcal H)|=p-2,
 \qquad
 \sigma_p=\frac{1-2/p}{(1-1/p)^2}
 =\frac{p(p-2)}{(p-1)^2}
 =1-\frac{1}{(p-1)^2}.
\]
Consequently
\[
 \mathfrak S(\{0,2\})
 =2\prod_{p>2}\left(1-\frac{1}{(p-1)^2}\right)>0.
\]
The positivity and convergence of this constant are local statements.  The
assertion that infinitely many integers $n$ make both $n$ and $n+2$ prime is
not a consequence of them.
\end{example}

\section{A finite divisor-sum criterion}

Let $\mathcal H$ be admissible, and let $W$ be a positive squarefree integer
divisible by every prime factor of $\Delta_{\mathcal H}$.  The finite product
map gives
\[
 |G_W(\mathcal H)|=\prod_{p\mid W}(p-\nu_p(\mathcal H))>0.
\]
Choose an integer $v_0=v_0(W)$ whose class lies in $G_W(\mathcal H)$.  Thus
$\gcd(v_0+h_i,W)=1$ for every $i$.  Let $N$ be a positive integer such that
\[
 I_N=\{n\in\mathbb Z:N\leq n<2N,\ n\equiv v_0\pmod W\}
\]
is nonempty and $n+h_i\geq2$ for $n\in I_N$ and all $i$.  Let
$\mathcal D\subset\mathbb N^k$ be finite, contain
$\boldsymbol 1=(1,\ldots,1)$, and satisfy
\[
 \mu^2(d_1\cdots d_k)=1,
 \qquad \gcd(d_1\cdots d_k,W)=1
 \quad(\boldsymbol d\in\mathcal D).
\]
The divisor-incidence map
\[
 T:\mathbb R^{\mathcal D}\longrightarrow\mathbb R^{I_N},
 \qquad
 (T\lambda)(n)=
 \sum_{\substack{\boldsymbol d\in\mathcal D\\
                  d_i\mid n+h_i\ (1\leq i\leq k)}}
 \lambda_{\boldsymbol d}
\]
is the finite algebra underlying the multidimensional Selberg weights in
\cite[(2.4), (4.2)--(4.3)]{Maynard}.  Put
\[
 R_{\mathcal H}(n)=\sum_{i=1}^k\mathbf 1_{\mathbb P}(n+h_i),
 \qquad (M_Rf)(n)=R_{\mathcal H}(n)f(n),
\]
where $\mathbb P$ is the set of positive primes, and define
\[
 A=T^*T,\qquad B=T^*M_RT,
\]
with respect to the standard Euclidean inner products.

\begin{proposition}[Finite divisor-sum Rayleigh principle]
For $\boldsymbol d,\boldsymbol e\in\mathcal D$, set
$q_i=\operatorname{lcm}(d_i,e_i)$.  The matrix entries are
\[
 A_{\boldsymbol d,\boldsymbol e}
 =\#\{n\in I_N:q_i\mid n+h_i\text{ for }1\leq i\leq k\}
\]
and, when the $q_i$ are pairwise coprime,
\[
 \left|A_{\boldsymbol d,\boldsymbol e}
       -\frac{N}{W\prod_{i=1}^kq_i}\right|\leq1.
\]
If they are not pairwise coprime, then
$A_{\boldsymbol d,\boldsymbol e}=0$.  Moreover,
\[
 B_{\boldsymbol d,\boldsymbol e}
 =\sum_{\substack{n\in I_N\\q_i\mid n+h_i\ (1\leq i\leq k)}}
 R_{\mathcal H}(n).
\]
For $\lambda\in\mathbb R^{\mathcal D}$, write
\[
 S_1(\lambda)=\|T\lambda\|_2^2,
 \qquad
 S_2(\lambda)=\langle T\lambda,M_RT\lambda\rangle.
\]
If $V=(\ker T)^\perp$ and $A_V,B_V$ denote the restrictions to $V$, then
$A_V$ is positive definite and
\[
 \max_{\lambda\notin\ker T}\frac{S_2(\lambda)}{S_1(\lambda)}
 =\lambda_{\max}\!\left(A_V^{-1/2}B_VA_V^{-1/2}\right)
 \leq \max_{n\in I_N}R_{\mathcal H}(n)\leq k.
\]
Consequently, if $m$ is a nonnegative integer and some $\lambda$ satisfies
\[
 S_1(\lambda)>0,
 \qquad S_2(\lambda)>mS_1(\lambda),
\]
then at least $m+1$ of $n+h_1,\ldots,n+h_k$ are prime for some $n\in I_N$.
\end{proposition}

\begin{proof}
The conditions on $\mathcal D$ imply that $W,q_1,\ldots,q_k$ are individually
squarefree and that $\gcd(W,q_i)=1$.  If a prime $p$ divides both $q_i$ and
$q_j$ with $i\ne j$, simultaneous divisibility would give
$h_i\equiv h_j\pmod p$.  But then $p\mid\Delta_{\mathcal H}$, hence $p\mid W$,
contrary to $\gcd(W,q_iq_j)=1$.  Thus the congruences counted by
$A_{\boldsymbol d,\boldsymbol e}$ are inconsistent unless the $q_i$ are
pairwise coprime.  In the remaining case, the Chinese remainder theorem gives
one residue class modulo $W\prod_iq_i$.  Its intersection with an interval of
$N$ consecutive integers has cardinality differing from
$N/(W\prod_iq_i)$ by at most $1$.  Expanding the two inner products gives the
displayed entries of $A$ and $B$.

Since $\boldsymbol 1\in\mathcal D$ and $I_N$ is nonempty, $T$ is nonzero.  If
$0\ne x\in V$, then
\[
 \langle x,A_Vx\rangle=\|Tx\|_2^2>0,
\]
so $A_V$ is positive definite.  Both $A$ and $B$ annihilate $\ker T$ and
preserve $V$.  Orthogonal projection onto $V$, followed by the substitution
$y=A_V^{1/2}x$, therefore reduces the generalized quotient to
\[
 \frac{\langle x,B_Vx\rangle}{\langle x,A_Vx\rangle}
 =\frac{\langle y,A_V^{-1/2}B_VA_V^{-1/2}y\rangle}
        {\langle y,y\rangle}.
\]
Rayleigh--Ritz gives the equality with the largest eigenvalue.  Finally,
\[
 S_2(\lambda)
 =\sum_{n\in I_N}R_{\mathcal H}(n)|(T\lambda)(n)|^2
 \leq\left(\max_{n\in I_N}R_{\mathcal H}(n)\right)S_1(\lambda).
\]
If $S_2(\lambda)>mS_1(\lambda)$, some term with
$(T\lambda)(n)\ne0$ has $R_{\mathcal H}(n)>m$.  This integer-valued prime
count is therefore at least $m+1$.
\end{proof}

Admissibility supplies the carrier $v_0$, but it supplies no strict inequality
between $S_2$ and $mS_1$.  Neither positivity of
$\mathfrak S(\mathcal H)$ nor the finite identities above implies the threshold
inequality.  No limiting operator or asymptotic estimate is asserted here.

\begin{thebibliography}{1}

\bibitem{Maynard}
J. Maynard,
\emph{Small gaps between primes},
Ann. of Math. (2) \textbf{181} (2015), no.~1, 383--413,
\href{https://doi.org/10.4007/annals.2015.181.1.7}
{doi:10.4007/annals.2015.181.1.7}.

\end{thebibliography}

\end{document}
